%!TEX root = TheoreticalPhysicsC.tex


Information about the organisation of the course, the \emph{coursework} and the \emph{examination}: be sure to look at Stud.ip.

We have a whole spectrum of topics to cover, namely
\begin{itemize}
\item Experimental findings
\item Postulates of quantum mechanics
\item Uncertainty
\item Schrödinger equation
\item simple potential problems
\item harmonic oscillator
\item hydrogen atom
\item identical particles
\item entanglement
\item Laws of thermodynamics
\item Carnot processes
\item Statistical description of equilibrium
\item Ideal (quantum) gases
\end{itemize}

I am not recommending a textbook to you: in my opinion the best book on quantum mechanics does not exist yet. There is no bad one, but equally no perfect one --- the authoritative one is this script. Since it is being written as the semester goes along, it contains mistakes; reports of any you find, sent by e-mail, will be rewarded at the end of the semester.

\section{Lecture: The Birth of Quantum Mechanics}
We find ourselves in a quantum universe. Until 1900, however, we did not notice the quantum aspects of the universe, because classical physics supplied us with a very good approximation to the physics of everyday life. Quantum mechanics would be far easier to understand if we could observe quantum aspects every day! (One of the greatest challenges in learning quantum mechanics is to understand why the approximation by classical physics works so well.) The keyword for that is \emph{decoherence}; it belongs to advanced quantum mechanics, which we shall not treat here. What stands in need of explanation is therefore not quantum mechanics, but the quality of the classical approximation.

Recently we have been able to watch a sensational development in physics: the first quantum computers have been built. For the first time we now have the chance to experiment directly with quantum systems (with the help of cloud services, e.g.\ the IBM Q experience). Quantum aspects will slowly become visible. Towards the end of the semester I would like to offer you a few exercises on real quantum computers.

Before we can turn to the latest results and their applications, let us first consider the development of QM and its formal structure.

\subsection{Prehistory}
First a remark on method, one that holds for the whole course. One can teach quantum mechanics by working through the wrong --- classical or semiclassical --- models first, in order to show afterwards that they do not work. I am no friend of that route: from the very beginning we look at everything correctly; nowhere do we work through the Bohr model of the atom. The historical experiments I shall name nonetheless, but only as a list.

The prehistory of quantum mechanics comprises a series of observations which gradually called into question the thesis that our universe is described by classical mechanics. Some of the most important of these contributions are
\begin{enumerate}
	\item \emph{Black-body radiation}. The spectrum of the emitted radiation cannot be explained with classical physics. The assumption of light quanta with energy $E=h\nu$ yields the right spectrum. Planck postulated the constant ``h'' (originally an ``auxiliary quantity'') in deriving the radiation law:
	\begin{equation*}
		\text{energy density of the radiation in the interval $d\nu$} = U(\nu)d\nu = \frac{8\pi h \nu^3}{c^3}\frac{d\nu}{e^{\frac{h\nu}{kT}}-1}
	\end{equation*}
	(1899). Today theoretical physics predominantly uses the \emph{reduced Planck constant}
	\begin{equation*}
		\eqbox{\hbar = \frac{h}{2\pi}}
	\end{equation*}
	instead.

	A black body reflects no incident radiation, yet at a temperature $T$ it radiates all the same, on account of its thermally excited atoms. The usual realisation is a box with a small hole: incident radiation vanishes inside through multiple reflection, and black-body radiation emerges from the hole. We shall derive this law ourselves in the statistical part of the course --- from the assumption that light is quantised.

	% fig_vl01_planck.tex — Planck gegen Rayleigh-Jeans: das Spektrum, an dem
% die klassische Physik scheitert, und der Grund dafür.
%
% Folgt den Konventionen von fig_vl11_kasten.tex.
%
% Was die Abbildung zeigt
%   (a) die Energiedichte U(nu) für drei Temperaturen, daneben die
%       klassische Vorhersage, die oben aus dem Bild herausläuft;
%   (b) warum: die mittlere Energie einer Mode, klassisch kT, tatsächlich
%       hnu/(exp(hnu/kT)-1) — oberhalb von hnu ~ kT friert sie ein.
%   Die beiden Panels hängen multiplikativ zusammen,
%       U(nu) = (8 pi nu^2/c^3) * <E>(nu,T),
%   also (a) = Modendichte x (b).  Genau der zweite Faktor ist es, den die
%   klassische Physik falsch bekommt; die Modendichte ist klassisch richtig.
%   Panel (b) ist zugleich das Verhältnis der beiden Kurven aus (a):
%       U/U_RJ = x/(e^x - 1)  mit  x = h nu/kT.
%
% Geometrie aus der Formel, nichts freihand
%   Planck (die Formel der Vorlesung):
%       U(nu) = (8 pi h nu^3/c^3) / (exp(h nu/kT) - 1)
%   Rayleigh-Jeans ist ihr Grenzfall h nu << kT:
%       U_RJ(nu) = 8 pi nu^2 k T / c^3
%   Mit nu = s * 1e14 Hz und U in 1e-15 J m^-3 Hz^-1 wird daraus
%       8 pi h/c^3 = 6.1806e-58  ->  U    = 0.61806 * s^3/(exp(a_T s) - 1),
%       8 pi k/c^3 = 1.28786e-47 ->  U_RJ = 1.28786e-4 * T * s^2,
%   mit a_T = h*1e14/(kT) = 4799.2/T, also a = 1.19980, 0.95984, 0.79987
%   für T = 4000, 5000, 6000 K.  Damit liegen die Maxima von selbst
%   richtig, wo Wien sie verlangt (h nu_max = 2.8214 kT):
%       s_max = 2.35, 2.94, 3.53  mit  U_max = 0.51, 0.99, 1.72,
%   d.h. die Lage wandert wie T und die Höhe wächst wie T^3.  Beides liest
%   man im Bild ab, ohne dass es hingeschrieben werden müsste.
%   Panel (b) ist dieselbe Planck-Formel in der dimensionslosen Variablen
%   x = h nu/kT:  <E>/kT = x/(exp(x)-1),  klassisch = 1.
%
% Entwurfsentscheidungen
%   * Farbe trägt die Temperatur, kalt -> heiß (spacecadet, cgblue,
%     birrenorange), und sonst nichts.  Beide klassischen Kurven —
%     Rayleigh-Jeans in (a) und die Gleichverteilungslinie <E> = kT in (b)
%     — sind neutral grau gestrichelt: gescheiterte Referenz, keine Daten.
%     Die eine Kurve in (b) ist temperaturunabhängig, also bekommt sie mit
%     birrengreen bewusst eine Farbe, die im Temperaturverlauf von (a)
%     nicht vorkommt.
%   * Rayleigh-Jeans ist für die *kälteste* der drei Temperaturen
%     gezeichnet, also für die flachste der drei klassischen Parabeln.
%     Sie verlässt den Rahmen trotzdem, bevor das zugehörige
%     Planck-Spektrum überhaupt sein Maximum erreicht, und übertrifft bei
%     nu = 2e14 Hz bereits das 6000-K-Spektrum: die Katastrophe hängt
%     nicht an der Temperatur.  Bei kleinen nu liegt die graue Kurve exakt
%     auf der navy Planck-Kurve — klassisch ist eben nicht alles falsch,
%     und genau das sagt (b) noch einmal quantitativ.
%   * Die Lücke zwischen kT und <E> in (b) ist schraffiert, nicht gefüllt
%     (Graustufendruck).  Sie ist als TikZ-Pfad gezeichnet und nicht mit
%     `fill between`: diese Bibliothek zerstört hier die Bounding Box der
%     tikzpicture, das rechte Panel wird dann abgeschnitten.
%   * Echte Einheiten (1e14 Hz, 1e-15 J m^-3 Hz^-1, Kelvin), damit die
%     Größenordnung mitkommt; 6000 K liegt bei der Sonnenoberfläche.
%   * Beschriftung direkt am Objekt statt in einer Legende.  Die
%     Achsenbeschriftungen stehen mittig unter bzw. neben der Achse statt
%     am Achsenende (tpcplot-Voreinstellung): am Achsenende würde die von
%     (a) in das rechte Panel hineinlaufen.

\begin{center}
\begin{tikzpicture}

% ─────────────────────────────────────────────────────────────
%  (a)  Das Spektrum: Planck für drei T, Rayleigh-Jeans daneben
% ─────────────────────────────────────────────────────────────
\begin{axis}[
    name=spectrum,
    tpcplot,
    width=8.2cm, height=5.8cm,
    xmin=0, xmax=8.8, ymin=0, ymax=2.6,
    xtick={0,2,4,6,8}, ytick={0,0.5,1,1.5,2,2.5},
    xlabel={$\nu\;[10^{14}\,\mathrm{Hz}]$},
    ylabel={$U(\nu)\;[10^{-15}\,\mathrm{J\,m^{-3}\,Hz^{-1}}]$},
    every axis x label/.style={at={(axis description cs:0.5,-0.13)},
                               anchor=north, font=\small},
    every axis y label/.style={at={(axis description cs:-0.135,0.5)},
                               rotate=90, anchor=south, font=\scriptsize},
    title={(a) \tpcifenglish{The spectrum}{Das Spektrum}},
]

  % Planck, U = 0.61806 s^3/(exp(a_T s) - 1)
  \addplot[spacecadet,   line width=1.1pt, domain=0.01:8.8, samples=220]
    {0.61806*x^3/(exp(1.199800*x)-1)};
  \addplot[cgblue,       line width=1.1pt, domain=0.01:8.8, samples=220]
    {0.61806*x^3/(exp(0.959840*x)-1)};
  \addplot[birrenorange, line width=1.1pt, domain=0.01:8.8, samples=220]
    {0.61806*x^3/(exp(0.799867*x)-1)};

  % Rayleigh-Jeans für T = 4000 K, U_RJ = 8 pi nu^2 kT/c^3: der Grenzfall
  % h nu << kT der Planck-Formel.  Grau, weil klassische Referenz.  Sie
  % wird zuletzt gezeichnet, damit man sie bis nu ~ 1e14 Hz genau auf der
  % navy Planck-Kurve liegen sieht, bevor sie sich löst und oben aus dem
  % Rahmen läuft.
  \addplot[birrengray, line width=0.9pt, densely dashed,
           domain=0:2.30, samples=60] {0.515144*x^2};

  % Temperaturen direkt an ihren Maxima (h nu_max = 2.8214 kT)
  \node[lbl, font=\scriptsize, text=birrenorange, anchor=south]
    at (axis cs:3.53,1.76) {$6000\,\mathrm{K}$};
  \node[lbl, font=\scriptsize, text=cgblue, anchor=south]
    at (axis cs:2.94,1.03) {$5000\,\mathrm{K}$};
  \node[lbl, font=\scriptsize, text=spacecadet, anchor=south]
    at (axis cs:2.35,0.55) {$4000\,\mathrm{K}$};

  % Die klassische Kurve und ihr Schicksal
  \node[lbl, font=\scriptsize, text=birrengray, anchor=north west,
        align=left, inner sep=1pt]
    at (axis cs:2.45,2.58)
    {Rayleigh--Jeans ($4000\,\mathrm{K}$)\\
     $U\propto\nu^{2}$: \tpcifenglish{UV catastrophe}{UV-Katastrophe}};
\end{axis}

% ─────────────────────────────────────────────────────────────
%  (b)  Warum: die mittlere Energie einer Mode
% ─────────────────────────────────────────────────────────────
\begin{axis}[
    tpcplot,
    at={(spectrum.south east)}, anchor=south west, xshift=10mm,
    width=5.4cm, height=5.8cm,
    xmin=0, xmax=6, ymin=0, ymax=1.28,
    xtick={0,1,2,3,4,5,6}, ytick={0,0.5,1},
    xlabel={$h\nu/kT$},
    ylabel={$\langle E\rangle/kT$},
    every axis x label/.style={at={(axis description cs:0.5,-0.13)},
                               anchor=north, font=\small},
    every axis y label/.style={at={(axis description cs:-0.16,0.5)},
                               rotate=90, anchor=south, font=\small},
    title={(b) \tpcifenglish{Energy per mode}{Energie pro Mode}},
]

  % Die Lücke zwischen klassischer Gleichverteilung und Planck:
  % schraffiert, nicht gefüllt.  Als Pfad, nicht mit `fill between`.
  \draw[forbidden, draw=none]
    (axis cs:0.02,1) -- plot[domain=0.02:6, samples=140, variable=\t]
      (axis cs:{\t},{\t/(exp(\t)-1)}) -- (axis cs:6,1) -- cycle;

  % Klassisch: jede Mode trägt kT, unabhängig von nu.
  \addplot[levellinegrey, domain=0:6, samples=2] {1};
  \node[lbl, font=\scriptsize, text=birrengray, anchor=south west,
        inner sep=1pt]
    at (axis cs:0.12,1.04)
    {\tpcifenglish{classical}{klassisch}: $\langle E\rangle=kT$};

  % Planck: <E>/kT = x/(exp(x)-1)
  \addplot[birrengreen, line width=1.2pt, domain=0.02:6, samples=200]
    {x/(exp(x)-1)};
  \node[lbl, font=\scriptsize, text=birrengreen, anchor=center,
        fill=white, fill opacity=0.92, text opacity=1, inner sep=1.5pt]
    at (axis cs:3.45,0.63)
    {$\langle E\rangle=\dfrac{h\nu}{e^{h\nu/kT}-1}$};

  % Was die Schraffur bedeutet
  \node[lbl, font=\scriptsize, text=birrengray, anchor=east,
        fill=white, fill opacity=0.92, text opacity=1, inner sep=1.5pt]
    at (axis cs:5.9,0.28)
    {\tpcifenglish{frozen out}{eingefroren}};
\end{axis}

\end{tikzpicture}
\end{center}


	Notice in~(a) how the grey classical curve at small $\nu$ still lies exactly on
	the $6000\,$K curve to which it belongs, and only then parts company with it;
	(b)~tells you where: as soon as $h\nu$ reaches the energy $kT$, the mode freezes out.

	\item \emph{Photoelectric effect}. The energy of the electrons released by light (photoelectrons) is determined by the frequency of the light alone (Lenard, 1899; Einstein 1905):
	\begin{equation*}
		\eqbox{E_{\text{el}} = \frac12 m v^2 = h\nu -W}
	\end{equation*}
	with binding energy $W$. This observation pointed to \emph{light quanta} with energy $E=h\nu$. Classically that is strange: everywhere else more intensity means more energy; here the energy depends on the frequency alone.
	\item \emph{Stability of atoms}, \emph{line spectra}, \emph{Franck-Hertz experiment} ($\sim$1911). In classical physics electron orbits in an atom would be unstable because of radiation losses. The radiation would be spread continuously over all frequencies. Line spectra and Franck-Hertz experiments indicate that only certain discrete energy levels are allowed in atoms. The reason for the instability is a result of electrodynamics that you have seen already: by the Liénard-Wiechert potentials, every accelerated point charge radiates. An electron on an orbit around the proton therefore loses energy all the time and ought to crash into the nucleus --- which it does not do. Only quantum mechanics can explain this; we shall discuss the line spectrum of hydrogen in detail later.
	\item \emph{Compton effect}. The scattering of light off electrons is described correctly by assuming light quanta with $E=h\nu$ and momentum $p=h/\lambda$, together with conservation of energy and momentum.
	\item \emph{Stern-Gerlach experiment} (1922). Angular momentum has only finitely many possible orientations relative to an applied magnetic field (space quantisation). This experiment matters particularly to us: precisely because angular momentum is quantised, we can build qubits today.
	\item \emph{Interference of electrons} (Davisson-Germer experiment, 1929; diffraction from crystals). This experiment gave an indication of the wave properties of the electron (postulated in 1924 by de Broglie).
\end{enumerate}

\begin{center}
\parbox[c][5cm][c]{14cm}{\begin{tcolorbox}[colback=banana,arc=0cm,frame hidden,enhanced]Matter and radiation possess both wave and particle properties.\end{tcolorbox}}
\end{center}



The year of birth of quantum mechanics is 1926. In that year Schrödinger published an exciting series of papers on wave mechanics and quantisation. A parallel approach began in 1925 with Heisenberg's reflections on the measurability of phase relations of particles in the Bohr model of the atom. That work inspired two teams to set up versions of quantum mechanics: on the one hand Born and Jordan with Heisenberg, leading to so-called matrix mechanics, and on the other hand Dirac. Later that same year Schrödinger shows the equivalence of the two approaches. The statistical interpretation of the theory is supplied by Born (likewise in 1926). With that the foundation of the theory is complete, and notions such as ``wave-particle duality'' have come into being. In my view Born's contribution is the most important aspect of the whole theory.

\begin{historical}
You are, by the way, in luck if you read German: the original papers of Planck, Schrödinger, Heisenberg and Born are written in German. I did not have that opportunity when I learnt quantum mechanics.
\end{historical}


\begin{center}
\parbox[c][5cm][c]{14cm}{\begin{tcolorbox}[colback=banana,arc=0cm,frame hidden,enhanced]Quantum mechanics is a \emph{probabilistic theory}. It therefore gives us predictions only in the form of \emph{probabilities}.\end{tcolorbox}}
\end{center}

This second statement stands above the first: ``wave and particle properties'' remains a little vague, whereas the probabilistic statement is precise and underlies every quantum phenomenon. Classically the predictions are definite --- the particle is at this place at this time --- and probabilities enter there purely through missing information. Please take this sentence away with you: in every oral examination I ask what the most important statement about quantum mechanics is.

Almost at the same time the theory took on the mathematical form it has today. Born, for one, took advice from Norbert Wiener on turning matrices into operators, so that ``aperiodic problems'' without a prescribed basis (e.g.\ a free particle) could be treated too. Hilbert gave a lecture course on quantum mechanics in Göttingen in the winter semester 1926/27. Alongside it his young assistant John von Neumann introduced the Hilbert space, density matrices and the entropy of quantum states, in a publication in the Göttinger Mathematische Nachrichten. Physicists did not read it at first. Wigner and Weyl clarified the role of group theory.

It is remarkable that, in spite of the theory's fantastic successes, its interpretation remained disputed. Despite many ideas and their development, the interpretation of QM is still controversial. Nevertheless QM is generally accepted today and gives us the best physical model we have in physics; to this day there is no contradiction with QM.

We now leave the exciting history of quantum mechanics behind. Our assumption is that quantum mechanics is the correct description of the universe.

Our next goal is to summarise the mathematical foundations of quantum mechanics. To reach that goal we must first treat \emph{probability} mathematically.

\subsection{Probability}
The numbers that quantum mechanics delivers as its results are \emph{probabilities}. How are such quantities to be understood, and how do we work with them? \emph{Probability theory} is the right framework for understanding random events and their interpretation. We shall build up that theory shortly. Do not expect anything conceptually new: I take an intuitive understanding for granted; what matters here is the \emph{notation} in which we shall later formulate quantum mechanics.


\subsubsection{Sample spaces, measures, moments}
In probability theory every system is assigned a set $X$, called the \emph{space of elementary events} or the \emph{sample space}.

\begin{example}
	Let us take the simplest case, in which $X$ is a finite set. For a die, for instance, the space of elementary events consists of the possible numbers of spots. Hence $X = \{1, 2, 3, 4, 5, 6\}$.
\end{example}

\begin{exercise}
	Describe the set $X$ to be assigned as the space of elementary events of \emph{two} dice.
\end{exercise}

A \emph{preparation} is characterised by giving a probability $p(x)$ for every $x\in X$.

\begin{example}
	For a fair die $p(x) = 1/6$ for all $x$. But there are other probabilities too, which then describe the preparations of \emph{loaded dice}.
\end{example}

The only restrictions on the probability distribution $p(x)$ are
\begin{equation}
	p(x) \ge 0 \quad \text{and}\quad \sum_{x\in X} p(x)  = 1.
\end{equation}

\begin{greybox}[Every experiment in three steps]
To the theorist \emph{all} experiments in physics look the same; they consist of three stages:
\begin{enumerate}
\item \emph{Preparation}: something is set up --- the die, or in the laboratory a beam of light, say. Preparations are probability functions $p$.
\item \emph{Execution}: the experiment runs. It maps probability functions onto probability functions.
\item \emph{Measurement}: at the end we look, and obtain one of the possible events. Measurements are functions $f$.
\end{enumerate}
An experimental physicist would not accept this simplification; in the abstract, though, every experiment looks like this. The first two steps can be gathered into \emph{one} more complicated preparation --- a loaded die is, after all, nothing but an ordinary die that somebody has done something to beforehand.
\end{greybox}

A \emph{measurement} (observable) is described by a function
\begin{equation}
	f:X\rightarrow \mathbb{R}
\end{equation}
on the sample space. Notice how generous that is: we impose no condition at all on a measurement; practically any function may count as one --- on the preparation, by contrast, two.
\begin{example}
	For a die the \emph{number of spots} $f(x) = x$ is a measurement.
\end{example}
``Is the number of spots even?'' is a measurement too. In probability theory such measurements are called \emph{random variables}. The value $f(x)$ is interpreted as the measured value obtained when one knows for certain that the elementary event $x$ has occurred. If this event occurs in $N$ measurements ($Np(x)$), the sum of the measured values comes out as $\sum_x N p(x) \cdot f(x)$. Divided by $N$ this gives us the formula for the expectation value:
\begin{equation}
	\text{expectation value} = \mathbb{E}[f] = \sum_{x\in X} p(x)f(x).
\end{equation}

An important role in quantum mechanics is played by the special case of random variables $f$ which can return only the results $0$ or $1$. Such measurements serve as a model for experiments, e.g.\ photon detectors, which measure the presence or absence of a particle. The detector says click or no click; quantum mechanics turns precisely on these two-valued measurements. These measurements are characterised by a set $S=\{x\in X\,|\, f(x) =1\}$. We often use the notation
\begin{equation}
	\chi_S(x) = \begin{cases}
		1, \quad x\in S, \\
		0, \quad \text{otherwise.}
	\end{cases}
\end{equation}
Such functions are called \emph{indicator functions}. The subsets $S\subset X$
are called events too (no longer elementary events). Note this: indicator functions \emph{are} measurements, not some separate object beside them. The set $S$ may be arbitrary; three limiting cases: $S = X$ gives the constant function $1$, $S = \emptyset$ the constant function $0$, and for a single elementary event $\chi_{\{j\}}(x) = \delta_{x,j}$ is the Kronecker delta.

The expectation value
\begin{equation}\label{eq:indicatorfunction}
	\mathbb{E}[\chi_S] = \sum_{x\in X} p(x)\chi_S(x) = \sum_{x\in S} p(x) \equiv \mathbb{P}[S]
\end{equation}
represents the probability that the event $x\in S$ takes place.

The function $\mathbb{P}$ has characteristic properties which follow at once from the definition. In the case where the subsets $S_j\subset X$ are disjoint, for instance,
\begin{equation}
	\mathbb{P}\left[\bigcup_j S_j \right] = \sum_{j} \mathbb{P}[S_j].
\end{equation}
Disjointness is essential: if the $S_j$ overlap, contributions get counted twice over and the equation becomes false. Mathematically speaking: $P$ is a measure, or a probability measure if one wants to emphasise the normalisation $\mathbb{P}[X] = 1$. We need no more measure theory than that; the notation is meant only to make the literature readable. Once the \emph{set function} $\mathbb{P}$ is known, the expectation-value functional $\mathbb{E}$ can be computed from it as well.

In the case that $X$ is a finite set, we can compute the expectation-value functional as follows
\begin{equation}
	\mathbb{E}[f] = \sum_{x\in X} \mathbb{P}[\{x\}] f(x).
\end{equation}
That is, $p(x) = \mathbb{P}[\{x\}]$.
\begin{exercise}
	Let $X = \{1,2,3\}$ and $p(x) = 1/3$ $\forall x\in X$. Compute the probability $\mathbb{P}[S]$ in the case $S=\{1,2\}$. Let $f(x) = x^2$. What is the expectation value
	\begin{equation}
		\mathbb{E}[x^2] = ?
	\end{equation}
\end{exercise}

So far $X$ has been finite. For quantum mechanics, though, $X = \mathbb{R}$ is the decisive case: a particle in one dimension has a position, and the possible positions are the real numbers; over an uncountable set $\sum_{x\in X}$ no longer has any meaning.

When $X = \mathbb{R}$, we use a different notation:
\begin{equation}
	\mathbb{E}[f] = \int \mathbb{P}(dx)\, f(x).
\end{equation}
(The sum is replaced by an integral.)
This notation is often written as
\begin{equation}\label{eq:ctsprob}
	\mathbb{E}[f] = \int p(x) f(x) \, dx
\end{equation}
as well, where $p(x)\,dx \equiv \mathbb{P}(dx)$ represents the probability distribution. I am not particularly fond of the first way of writing it and shall hardly use it; you should merely be able to recognise it.

\begin{example}
	The probability distribution
	\begin{equation}
		p(x) \equiv \frac{1}{\sqrt{2\pi}}e^{-\frac{1}{2}x^2}
	\end{equation}
	represents a \emph{normal} or \emph{Gaussian distribution}. The probability that $x\in S$, where $S= \{x\,|\, x\ge 0\}$, is computed by the following integral (by Eq.~\ref{eq:indicatorfunction} and Eq.~\ref{eq:ctsprob}):
	\begin{equation}
		\begin{split}
			\mathbb{P}[S] = \mathbb{E}[\chi_S] &= \frac{1}{\sqrt{2\pi}}\int_{-\infty}^\infty \chi_S(x) e^{-\frac{1}{2}x^2}\, dx \\ &= \frac{1}{\sqrt{2\pi}}\int_{-\infty}^0 \chi_S(x) e^{-\frac{1}{2}x^2}\, dx + \frac{1}{\sqrt{2\pi}}\int_{0}^\infty \chi_S(x) e^{-\frac{1}{2}x^2}\, dx = \frac{1}{2}.
		\end{split}
	\end{equation}
\end{example}

Intuitively the result was clear in advance: one half, because the Gaussian distribution is symmetric. Be careful with that argument, though --- it works \emph{only} for symmetric densities. Put physically: is the particle to the left or to the right of the origin?

For measures on $\mathbb{R}$ it is often particularly easy to compute the expectation values of the coordinate function and its powers. These are called \emph{moments}. For a probability measure $\mathbb{P}(dx) = p(x)dx$ on $\mathbb{R}$ one calls
\begin{equation}
	M_n = \int x^n p(x)\, dx
\end{equation}
the \emph{$n$-th moment} of $\mathbb{P}$. One always has $M_0 = 1$, because of the normalisation to total probability 1. (Warning: none of the other moments need exist, and some may give divergent integrals.) Geometrically $M_0$ is the area under the curve.

The moments that matter most are $M_1$ (the \emph{mean}) and $M_2$. In practice the second moment is often replaced by the \emph{variance}:
\begin{equation}
	\text{Var}[\mathbb{P}] \equiv \mathbb{E}[(x-M_1)^2] \equiv \int (x^2-2M_1 x +M_1^2) p(x)\, dx \underset{\text{Exercise!}}{=} M_2-M_1^2.
\end{equation}
From the definition it is clear that variances are always $\ge 0$ and vanish only when $\mathbb{P}$ is concentrated on a single point.

In quantum mechanics the variance plays a larger role than the second moment; we shall meet it again with the uncertainty relation. Mean and variance carry the same information as $M_1$ and $M_2$.
